\chapter{Several variables: local maps and conservation}
\label{learn:multivariable}

\begin{learningcard}{Prerequisites and question}
Use derivatives and integrals from Chapter~\ref{learn:calculus}, and matrix
maps and dot products from Chapter~\ref{learn:linear_algebra}. A temperature
field depends on position, while a moving fluid carries a direction at every
position. How can local rates describe a changing quantity spread across a
region? We work first in Euclidean space with smooth functions and regular
boundaries; singular fields require additional analysis.
\end{learningcard}

\section{A partial derivative controls one input}
A scalar field assigns one real value to each point in a region of $\RR^n$.
For $f(x,y)$, the partial derivative with respect to $x$ is
\[
 f_x(a,b)=\lim_{h\to0}\frac{f(a+h,b)-f(a,b)}h,
\]
when the limit exists. The second input stays fixed. The derivative $f_y$
is defined by varying only the second input. For
$f(x,y)=x^2+3xy+y^2$, ordinary differentiation gives
$f_x=2x+3y$ and $f_y=3x+2y$.

The gradient is the vector $\nabla f=(f_x,f_y)$ in Cartesian coordinates.
If $f$ is differentiable at a point $p$, then
\[
 f(p+h)=f(p)+\nabla f(p)\cdot h+r(h),\qquad
 \frac{|r(h)|}{\|h\|}\longrightarrow0\quad(h\to0).
\]
This condition controls all directions together. The existence of partial
derivatives at one point alone provides weaker information. For example,
set $g(x,y)=xy/(x^2+y^2)$ away from the origin and $g(0,0)=0$.
Both coordinate-axis restrictions vanish, so both partials at the origin
are zero. Along $y=x\ne0$, however, $g=1/2$, so $g$ is discontinuous
at the origin and therefore fails to be differentiable. Continuous first
partials in a neighborhood are a sufficient condition for differentiability.

\section{Worked example: choose a direction of change}
For the polynomial above, $\nabla f(1,0)=(2,3)$. Along a unit direction
$u$, the directional derivative is $\nabla f(1,0)\cdot u$.
The Cauchy--Schwarz inequality bounds its magnitude by $\sqrt{13}$.
For a unit vector $u$, the needed bound follows directly from
$0\leq\|v-(v\cdot u)u\|^2=\|v\|^2-(v\cdot u)^2$;
therefore $|v\cdot u|\leq\|v\|$.
The largest increase occurs at $u=(2,3)/\sqrt{13}$ and equals
$\sqrt{13}$. The direction $(3,-2)/\sqrt{13}$ has zero first-order
change because its dot product with the gradient vanishes.

For an actual displacement $h=(0.01,-0.02)$, the linear prediction is
$2(0.01)+3(-0.02)=-0.04$. Expanding the polynomial shows an exact
additional term $h_x^2+3h_xh_y+h_y^2=-0.0001$, so
$f(1.01,-0.02)-f(1,0)=-0.0401$. The calculation measures the error
instead of treating a local approximation as an identity.

\section{The Jacobian organizes several outputs}
For $F\colon\RR^n\to\RR^m$, its Jacobian matrix has entries
$(DF)_{ij}=\partial F_i/\partial x_j$. Differentiability means
$F(p+h)=F(p)+DF(p)h+r(h)$ with $\|r(h)\|/\|h\|\to0$.
The chain rule becomes matrix multiplication:
\[
 D(G\circ F)(p)=DG(F(p))\,DF(p).
\]
The dimensions enforce the order: an $n$-component displacement enters
$DF$ first, producing an $m$-component displacement for $DG$.

Polar coordinates provide a concrete example:
$F(r,\theta)=(r\cos\theta,r\sin\theta)$. Then
\[
 DF=\begin{pmatrix}\cos\theta&-r\sin\theta\\
 \sin\theta&r\cos\theta\end{pmatrix},\qquad \det DF=r.
\]
The determinant measures signed local area scaling; its absolute value
enters an area integral. Away from $r=0$, a small coordinate rectangle
with sides $dr,d\theta$ therefore has area approximately $r\,dr\,d\theta$.
Coordinate degeneracy at the origin explains why a single global inverse
polar chart requires care.

\begin{learningcard}{Historical situation: notation as shared equipment}
Edwin Bidwell Wilson's 1901 \emph{Vector Analysis}, based on J. Willard
Gibbs's lectures, explicitly presents a textbook for students of mathematics
and physics. Its chapters on differentiation and integration show vector
methods being organized for instruction. The title page documents the
collaboration and pedagogical setting; the chapter organization supplies
evidence of a teachable framework rather than a claim of sole invention.
\footnote{E. B. Wilson, \emph{Vector Analysis: A Text-Book for the Use of
Students of Mathematics and Physics, Founded upon the Lectures of
J. Willard Gibbs} (1901), title page and chapters on differentiation and
integration; digitized text:
\url{https://archive.org/details/vectoranalysiste00gibbiala}.}
Gibbs's preface describes a privately circulated pamphlet printed in
1881--1884 and requests for a fuller public treatment. He then identifies
Wilson as a listener to his 1899--1900 course who undertook the textbook.
Wilson's preface explains his decision to distribute applications through
the exposition as the necessary mathematics became available. These
documented choices show how lectures, limited circulation, and student
needs shaped the published teaching form. Our Jacobian notation and
examples give a modern route through that subject.
\end{learningcard}

\section{Worked example: integration needs the area factor}
The double integral $\iint_D f\,dA$ accumulates values times small areas
and is defined by a limiting sum under suitable integrability assumptions.
For a continuous function on a rectangle, iterated one-variable integration
computes the double integral. Changing variables by a continuously
differentiable one-to-one map with a nonzero Jacobian determinant on the
interior introduces $|\det DF|$.

Compute the integral of $x^2+y^2$ over the unit disk. Polar coordinates
cover the interior apart from the origin and a seam of area zero, which
leave the integral unchanged. The result is
\[
 \iint_{x^2+y^2\leq1}(x^2+y^2)\,dA
 =\int_0^{2\pi}\int_0^1r^2r\,dr\,d\theta
 =2\pi\left[\frac{r^4}{4}\right]_0^1=\frac\pi2.
\]
Dividing by the disk's area $\pi$ gives the mean squared distance from
the center, $1/2$. Omitting the Jacobian factor would count equal radial
widths as equal areas despite their different circumferences.

\section{Divergence turns flux into a local balance}
A vector field $J=(J_1,J_2,J_3)$ assigns a vector to each point. Its
divergence is $\nabla\cdot J=\partial_xJ_1+\partial_yJ_2+\partial_zJ_3$.
For a continuously differentiable field on a neighborhood of a bounded region
$V$ with piecewise smooth boundary, the divergence theorem states
\[
 \int_V\nabla\cdot J\,dV=\int_{\partial V}J\cdot n\,dS,
\]
where $n$ is the outward unit normal and $dS$ is surface area. The right
side is outward flux. On a small rectangular box, opposite-face differences
in each component produce the corresponding partial derivative times volume;
summing boxes cancels shared interior faces. This cancellation explains the
theorem's structure; taking the limit rigorously requires the stated regularity.

Let $\rho(x,t)$ be density, $J(x,t)$ its outward transport flux, and
$s(x,t)$ a local production rate per volume. For a fixed region,
\[
 \frac{d}{dt}\int_V\rho\,dV=-\int_{\partial V}J\cdot n\,dS+\int_Vs\,dV.
\]
Assuming sufficient smoothness to move the derivative inside the integral
and applying the divergence theorem gives
$\int_V(\partial_t\rho+\nabla\cdot J-s)\,dV=0$.
If the continuous integrand integrates to zero on every sufficiently small
region, it vanishes pointwise; otherwise a neighborhood of a strictly
positive or negative value would have a nonzero integral. Hence
$\partial_t\rho+\nabla\cdot J=s$. With $s=0$ and zero boundary flux,
total mass is constant. Chapter~\ref{learn:dynamics} uses this balance to
derive a heat equation and a numerical conservation check.

\section{Exercises with full solutions}
\paragraph{1. Differentiate along a trajectory.}
For $f(x,y)=x^2+y^2$ and $(x(t),y(t))=(\cos t,\sin t)$, find $df/dt$.

\paragraph{Solution.}
The chain rule gives $df/dt=2x x'+2y y'=2\cos t(-\sin t)+
2\sin t\cos t=0$. Direct substitution gives $f=1$, confirming the
result independently. The gradient is radial and the trajectory's velocity
is tangent to the circle, so their dot product vanishes.

\paragraph{2. Verify a divergence theorem example.}
For $J=(x,y,z)$ on the unit cube $[0,1]^3$, compute both sides.

\paragraph{Solution.}
The divergence is three, so its volume integral is three. On each face
$x=0$, $y=0$, or $z=0$, the normal component is zero. On each opposite
face the outward normal component is one and the area is one. The total
outward flux is therefore $1+1+1=3$, matching the volume integral.

\paragraph{3. Identify the conservation boundary.}
In a one-dimensional interval $[0,L]$, suppose $\rho_t+J_x=0$.
Express the mass derivative and identify a sufficient boundary condition.

\paragraph{Solution.}
Integrating and using the fundamental theorem yields
$M'(t)=-[J]_0^L=J(0,t)-J(L,t)$, where $M=\int_0^L\rho\,dx$.
Equal endpoint fluxes guarantee $M'=0$; zero flux at both endpoints is one
such condition. Periodic flux is another. The local equation alone allows
mass to leave through the boundary, so the boundary convention belongs in
any global conservation claim.
